\section{Introducción}
\label{ch:introduccion}

En el presente documento se presenta el proyecto  \textbf{\NombreProyecto}, un sistema de chat automatizado con inteligencia artificial diseñado para facilitar la cotización de productos a través de distintos canales digitales, así como la gestión administrativa de empresas y clientes. El sistema integra tecnologías de backend, frontend, base de datos e inteligencia artificial, ofreciendo un panel administrativo que consolida la información de múltiples empresas y el historial de cotizaciones generadas.

El objetivo principal del proyecto es reducir el tiempo de atención al cliente en el proceso de cotización de productos, mejorando así la experiencia del usuario final mediante respuestas más rápidas y precisas. Esta propuesta resulta especialmente atractiva para empresas que aún no cuentan con un sistema de chat automatizado y buscan optimizar su servicio al cliente.

A lo largo de este documento se definen todos los procesos e información relevante del proyecto, abordando desde el problema que se pretende resolver, los objetivos establecidos y la justificación de la propuesta, hasta los alcances y limitaciones que delimitan su desarrollo.

\subsection{Objetivos}
\label{sec:objetivos}

Los objetivos son las metas que se esperan alcanzar con la realización de la investigación. Deben ser claros, medibles, alcanzables, relevantes y con un tiempo definido (SMART, por sus siglas en inglés). Se dividen en un objetivo general y varios objetivos específicos.

\subsubsection{Objetivo General}
\label{subsec:objetivo_general}

El objetivo general expresa el propósito fundamental del estudio en términos amplios. Es la meta principal que se busca lograr y generalmente abarca la totalidad de la investigación. Debe ser coherente con la definición del problema.
\begin{itemize}
    \item \textbf{Objetivo:} Desarrollar un sistema cotizador de productos a través de Inteligencia Artificial que automatice el proceso de consulta, análisis y generación de cotizaciones, permitiendo mejorar la rapidez, precisión y atención al usuario, y teniendo como entregable un sistema funcional evaluado.
\end{itemize}

\subsubsection{Objetivos Específicos}
\label{subsec:objetivos_especificos}

Los objetivos específicos detallan los pasos o las acciones concretas que se llevarán a cabo para alcanzar el objetivo general. Son tareas más pequeñas y desglosadas que, al ser cumplidas, permiten la consecución del objetivo principal.
\begin{itemize}
    \item \textbf{Objetivo 1:} Analizar los requerimientos funcionales y tecnológicos del sistema, generando como entregable un documento de especificación de requisitos.
    \item \textbf{Objetivo 2:} Diseñar y desarrollar el sistema cotizador con Inteligencia Artificial, generando como entregable un prototipo funcional con sus módulos principales de operación.
    \item \textbf{Objetivo 3:} Validar el desempeño del sistema mediante pruebas de eficiencia, precisión y satisfacción del usuario, generando como entregable un informe de resultados y propuestas de mejora.
\end{itemize}

\subsection{Planteamiento del Problema}
\label{sec:problema}

El problema que se busca resolver es la ineficiencia en el proceso de cotización de productos en empresas medianas y pequeñas. En la actualidad, dicho proceso se realiza principalmente a través de mensajería instantánea, sin una correcta automatización. Son los propios trabajadores de las empresas quienes atienden manualmente las solicitudes de cotización, lo que genera retardos significativos en el intervalo entre la solicitud del cliente y la respuesta. Esta situación afecta directamente diversos aspectos de la atención al cliente, como la velocidad de respuesta, la disponibilidad fuera del horario laboral y la consistencia de la información proporcionada.

Frente a esta problemática, la presente investigación propone automatizar el proceso de cotización mediante un sistema de respuesta automatizada basado en inteligencia artificial, complementado con un panel administrativo que permita a las empresas revisar y modificar la información relacionada con sus propios productos y las cotizaciones realizadas por los clientes.

Con base en lo anterior, la pregunta de investigación que guía este trabajo es: ¿De qué manera la implementación de un sistema de chat automatizado con inteligencia artificial puede reducir el tiempo de respuesta en el proceso de cotización de productos en empresas de mediano y pequeño tamaño?

\subsection{Antecedentes}
\label{sec:antecedentes}

En esta sección, se presenta una revisión concisa y pertinente de investigaciones previas, literatura existente y el estado del arte relacionado con el tema de la tesis. El objetivo es situar el trabajo actual dentro del conocimiento acumulado, identificar vacíos en la investigación y demostrar la originalidad y la necesidad del estudio propuesto. Se deben citar adecuadamente las fuentes, mostrando cómo el trabajo se construye sobre cimientos previos o cómo aborda aspectos no resueltos.

\subsection{Antecedentes de la empresa}
\label{sec:antecedentes_de_la_empresa}

\textbf{\UnidadEconomica} es una empresa enfocada en brindar servicios de soporte técnico y formación profesional en el área de tecnologías de la información. Su trabajo principal consiste en resolver problemas operativos de hardware y software para empresas y particulares, además de ofrecer cursos prácticos para mejorar el uso de herramientas digitales.

Su misión es impulsar la productividad de sus clientes mediante el dominio de la tecnología, buscando siempre que las personas y negocios sean más eficientes al utilizar herramientas digitales. Para lograr esto, se enfocan en cerrar la brecha técnica a través de la capacitación continua y la certificación oficial de competencias.

En cuanto a sus actividades, la empresa combina el mantenimiento técnico con la enseñanza. Por un lado, se encargan de reparar equipos, optimizar sistemas y proteger la infraestructura informática de sus clientes. Por otro lado, operan como un centro de capacitación y certificación donde enseñan a manejar programas de oficina y diseño, validando profesionalmente los conocimientos adquiridos por sus alumnos.

Actualmente me encuentro trabajando en uno de sus proyectos que consiste en el \textbf{\NombreProyecto}, siendo esta una de las soluciones tecnológicas que ofrecen a distintas empresas y/o negocios.

\subsection{Antecedentes teóricos}
\label{sec:antecedentes_teoricos}

\subsubsection{Inteligencia Artificial y Procesamiento de Lenguaje Natural}
\label{subsec:inteligencia_artificial}

\subsubsection{Arquitectura RAG (Retrieval-Augmented Generation)}
\label{subsec:arquitectura_rag}

\subsubsection{WhatsApp Business Platform}
\label{subsec:whatsapp_business_platform}

\subsection{Tabla comparativa}
\label{sec:tabla_comparativa}

Con el fin de evidenciar las mejoras que introduce el presente proyecto, en la Tabla \ref{tab:comparativa} se contrastan los procesos realizados de manera tradicional (sin automatización) frente a la operación del sistema propuesto, abarcando aspectos como el tiempo de respuesta, disponibilidad, registro de conversaciones y gestión de productos.

\begin{table}[h]
\centering
\caption{Comparativa entre el proceso tradicional y el sistema propuesto}
\label{tab:comparativa}
\begin{tabular}{|p{4cm}|p{5cm}|p{5cm}|}
\hline
\textbf{Aspecto} & \textbf{Antes (sin el sistema)} & \textbf{Después (con el sistema)} \\
\hline
Proceso de cotización & Realizado manualmente por empleados vía WhatsApp o llamada telefónica & Proceso automatizado mediante chatbot con IA vía WhatsApp/widget \\
\hline
Tiempo de respuesta al cliente & Varía desde minutos hasta horas & Inmediato o en algunos segundos \\
\hline
Disponibilidad del servicio & Únicamente en horario laboral & Disponible las 24 horas de los 7 días de la semana \\
\hline
Registro de conversaciones & Nulo o bastante informal (capturas de pantalla) & Registro centralizado en base de datos \\
\hline
Cotizaciones generadas & Manualmente, propensas a tener errores de cálculo & Automáticamente, con cálculos precisos \\
\hline
Almacenamiento de cotizaciones & A papel, en archivos sueltos o mensajes de WhatsApp & En base de datos, con posibilidad de consultar y exportar a PDF \\
\hline
Gestión de productos & En Excel o a papel & Panel administrativo con CRUD \\
\hline
Disponibilidad de información & Distintas maneras de manejar la distinta información & Panel centralizado con la información necesaria \\
\hline
\end{tabular}
\end{table}

\subsection{Justificación}
\label{sec:justificacion}

El presente proyecto se justifica desde distintas perspectivas por los beneficios que puede aportar al automatizar el proceso de cotización de productos en empresas medianas y pequeñas.

Desde el punto de vista tecnológico, el proyecto representa la integración de una tecnología emergente como lo es la inteligencia artificial en un sector productivo como lo son las empresas o negocios. La integración de un chatbot basado en IA se encarga de modernizar los canales de atención al cliente, ofreciendo un servicio más rápido y eficiente.

Desde el punto de vista económico, la automatización del proceso de cotización reduce los intervalos de tiempo entre la solicitud del cliente y la respuesta, lo que puede traducirse en una optimización de los empleados de la empresa, liberándolos de tareas repetitivas.

Desde el punto de vista práctico, el sistema genera un reporte que integra los productos cotizados y la información necesaria sobre estos, además de almacenar toda la información sobre los productos de cierto negocio, facilitando el seguimiento de las cotizaciones realizadas y la disponibilidad de los productos.

Debido a lo anterior, invertir tiempo y recursos en este proyecto vale la pena porque ofrece una solución concreta a una situación real que enfrentan empresas o negocios en la actualidad.

\subsection{Alcances y Limitaciones}
\label{sec:alcances_limitaciones}

\subsubsection{Alcances}
\label{subsec:alcances}

\paragraph{Sobre el sistema en general}
\begin{itemize}
    \item El sistema consistirá en un chatbot con inteligencia artificial para atención al cliente, consulta de productos y generación automática de cotizaciones.
    \item La plataforma será multiempresa, permitiendo que múltiples negocios operen de manera independiente dentro del mismo sistema.
    \item El sistema incluirá un panel administrativo web para la gestión de empresas, usuarios, productos, conversaciones y reportes.
\end{itemize}
\paragraph{Sobre los canales de comunicación}
\begin{itemize}
    \item El chatbot estará disponible a través de un widget web embebible en los sitios de las empresas.
    \item El chatbot estará disponible a través del WhatsApp de cada empresa.
    \item Ambos canales de comunicación (WhatsApp y widget web) compartirán el mismo motor de IA y registrarán las conversaciones en la base de datos.
\end{itemize}
\paragraph{Sobre los perfiles de usuario}
\begin{itemize}
    \item El sistema contará con tres perfiles internos: administrador general, administrador de empresa y usuario de empresa.
    \item Los clientes finales interactuarán únicamente a través del widget web o WhatsApp, sin acceso al panel administrativo.
    \item El administrador general tendrá visibilidad y control sobre todas las empresas del sistema.
    \item El administrador de empresa podrá gestionar únicamente su propia empresa, sus productos y sus usuarios.
    \item El usuario de empresa podrá gestionar productos y visualizar reportes, con permisos restringidos.
\end{itemize}
\paragraph{Sobre la gestión de productos}
\begin{itemize}
    \item El sistema permitirá el registro manual de productos mediante formulario.
    \item El sistema permitirá la importación masiva de productos mediante archivos CSV.
    \item El sistema permitirá la importación asistida con IA para archivos sin formato predefinido.
    \item Se podrán editar, activar, desactivar o eliminar productos.
    \item Se podrán gestionar categorías, precios, descripciones, imágenes, stock y unidad de medida.
\end{itemize}
\paragraph{Sobre las conversaciones y cotizaciones}
\begin{itemize}
    \item El chatbot responderá exclusivamente con base en el catálogo de productos real de cada empresa.
    \item El sistema generará cotizaciones automáticas a partir de los productos seleccionados por el cliente durante la conversación.
    \item Las cotizaciones incluirán cálculo de cantidades, precios, subtotales y totales.
    \item Las cotizaciones podrán ser guardadas, consultadas en el historial y exportadas a PDF.
    \item Se registrará el historial completo de mensajes de cada conversación.
\end{itemize}
\paragraph{Sobre los reportes y auditoría}
\begin{itemize}
    \item El sistema generará reportes de uso por empresa: tokens consumidos, conversaciones atendidas, productos consultados y cotizaciones generadas.
    \item Los reportes podrán filtrarse por rango de fechas.
    \item Los reportes podrán exportarse a Excel y PDF.
    \item Se registrarán logs de auditoría incluyendo fecha, hora, usuario o empresa asociada, y resultado de la acción.
\end{itemize}
\paragraph{Sobre el widget web}
\begin{itemize}
    \item Se generará un script embebible por empresa para integrar el widget en sitios web externos.   
    \item El widget permitirá configuración visual personalizada: color principal, logo, nombre del asistente, mensaje inicial, posición en pantalla, estilo de burbuja, mensajes de error y carga.
    \item Se implementará protección contra abuso o llamadas excesivas al widget.
\end{itemize}
\paragraph{Fronteras}
\begin{itemize}
    \item No se incluyen otros canales de comunicación como Messenger, Telegram, SMS o correo electronico.
    \item No se soportan imagenes, videos, audios ni documentos en el chat, unicamente texto.
    \item No se contempla app móvil, el acceso es via web para el panel administrativo y via widget/WhatsApp para el chatbot.
    \item El sistema unicamente operará en español, no se contemplan otros idiomas.
\end{itemize}

\subsubsection{Limitaciones}
\label{subsec:limitaciones}

\begin{itemize}
    \item El desarrollo del sistema debe completarse dentro del período establecido por el calendario académico (Mayo-Agosto 2026), lo que puede restringir la profundidad de ciertos módulos o la implementación de funcionalidades adicionales no críticas.
    \item La calidad de las respuestas del chatbot depende de la calidad y completitud del catálogo de productos ingresado por cada empresa. Datos incompletos o inconsistentes pueden afectar el rendimiento de la IA.
    \item El despliegue inicial contempla un solo Droplet de Digital Ocean, lo que podría limitar la capacidad de respuesta ante picos de demanda superiores a los 200 conversaciones simultáneas estimadas como base.
\end{itemize}

\subsection{Temas tentativos del marco teórico}
\label{sec:temas_tentativos_marco_teorico}

% \subsection{Ejemplos de Tablas e Inclusión de Imágenes}
% \label{sec:ejemplos_elementos}

% Para ilustrar el uso de elementos visuales en su tesis, a continuación se presentan ejemplos básicos de cómo incluir tablas e imágenes en LaTeX.

% \subsubsection{Inclusión de Tablas}

% Las tablas son herramientas fundamentales para presentar datos de manera organizada y comprensible.
% \begin{table}[H]
%     \centering
%     \caption{Clasificación de Sensores Comunes en Ingeniería.}
%     \label{tab:sensores_comunes}
%     \begin{tabular}{|l|c|c|}
%         \hline
%         \textbf{Tipo de Sensor} & \textbf{Magnitud Medida} & \textbf{Principio de Funcionamiento} \\
%         \hline
%         Termopar & Temperatura & Efecto Seebeck \\
%         Extensómetro & Deformación & Resistencia eléctrica \\
%         Acelerómetro & Aceleración & Efecto piezoeléctrico \\
%         Encoder rotatorio & Posición angular & Óptico/Magnético \\
%         \hline
%     \end{tabular}
%     \caption*{Nota: Esta tabla es un ejemplo para demostrar la estructura básica de una tabla en LaTeX.}
% \end{table}
% Como se puede observar en la Tabla \ref{tab:sensores_comunes}, la información se presenta de forma clara y estructurada. Es importante añadir un \verb|`\caption`|  para describir la tabla y un  \verb|`\label`|  para referenciarla en el texto.

% \subsubsection{Inclusión de Imágenes}

% Las imágenes, gráficos o diagramas son cruciales para complementar la explicación textual y facilitar la comprensión de conceptos complejos o resultados visuales.
% \begin{figure}[H]
%     \centering
%     \includegraphics[width=0.7\textwidth]{img/diagrama_proceso}
%     \caption{Diagrama de un proceso de control automático.}
%     \label{fig:diagrama_control}
% \end{figure}
% % La Figura \ref{fig:diagrama_control} ilustra un ejemplo de un diagrama de un proceso de control. Para incluir una imagen, se utiliza el entorno `figure` y el comando `\includegraphics`. Es fundamental especificar la ruta de la imagen (en este caso, `images/diagrama_proceso.png`) y un `width` para controlar su tamaño. Al igual que con las tablas, un `\caption` y un `\label` son indispensables. Es crucial que cada imagen tenga una calidad adecuada y que sea relevante para el texto.

% \begin{figure}[H]
%     \centering
%     \begin{subfigure}[b]{0.45\textwidth}
%         \includegraphics[width=\textwidth]{img/grafico_datos1}
%         \caption{Datos de Temperatura.}
%         \label{fig:subfig_temp}
%     \end{subfigure}
%     \hfill
%     \begin{subfigure}[b]{0.45\textwidth}
%         \includegraphics[width=\textwidth]{img/grafico_datos2}
%         \caption{Datos de Presión.}
%         \label{fig:subfig_pres}
%     \end{subfigure}
%     \caption{Comparación de datos operativos de un sistema (Temperatura vs. Presión).}
%     \label{fig:comparacion_datos}
% \end{figure}
% Incluso es posible combinar varias imágenes en una sola figura, como se muestra en la Figura \ref{fig:comparacion_datos}, utilizando el paquete `subcaption`.

% Asegúrense de que todas las figuras y tablas estén debidamente referenciadas en el texto y que su contenido sea autoexplicativo, utilizando los pies de figura y títulos de tabla de manera efectiva.